\chapter{LITERATURE REVIEW}

\section{General}

The literature review chapter demonstrates a thorough knowledge of the interdisciplinary areas fundamental to this thesis, specifically coastal hydrodynamic modelling, deep learning, and physics-informed surrogate frameworks. The primary objective is to critically evaluate the progression from traditional physics-based numerical solvers to contemporary machine learning approaches in fluid mechanics, highlighting the strengths and inherent limitations of each paradigm. By synthesizing recent advancements in Physics-Informed Neural Networks (PINNs), Graph Neural Networks (GNNs), and spectral bias mitigation techniques, this chapter provides the foundational arguments to support the study's focus: the development of a fully parametric Physics-Informed Graph Neural Network (PI-GNN) for estuarine tidal hydrodynamics. 

\section{Historical Background}

\subsection{Traditional Numerical Hydrodynamic Modelling}

Historically, the prediction of estuarine and coastal hydrodynamics has relied exclusively on physics-based numerical models. These models fundamentally solve the Navier--Stokes equations, most commonly simplified to the depth-averaged Shallow Water Equations (SWE) for free-surface flows where horizontal scales vastly exceed the vertical scale \cite{Leveque2002}. Industry-standard solvers such as Delft3D Flexible Mesh (Delft3D FM) \cite{Delft3DFM2011} and TELEMAC-2D \cite{Hervouet2000} employ advanced numerical schemes, such as the Finite Volume Method (FVM), to discretise the spatial domain into highly irregular, unstructured meshes. These unstructured grids allow for high-resolution representation of complex coastlines, deep navigational channels, and highly variable bathymetry, which are critical for capturing non-linear tidal propagation and riverine interactions. 

When properly calibrated against field measurements, these numerical models exhibit remarkable accuracy and physical fidelity \cite{Bates2010, Quirogaetal2020}. However, their profound limitation lies in their computational expense. The necessity to satisfy strict stability criteria (e.g., the Courant--Friedrichs--Lewy condition) forces these models to operate on very small time steps, often requiring hours or days of high-performance computing resources to simulate a single tidal cycle. As the global imperative to assess climate-driven hydrological extremes grows---driven by rising sea levels and intensified storm surges \cite{IPCC2021, Vousdoukas2018, Alfieri2017}---the prohibitive cost of these traditional models renders them fundamentally unsuited for large-scale probabilistic risk assessments, real-time flood warning systems, and rapid scenario exploration.

\subsection{The Rise of Data-Driven Machine Learning in Hydrology}

To circumvent the computational bottlenecks of traditional solvers, the past decade has seen a paradigm shift towards data-driven surrogate modelling, spurred by broader advancements in deep learning \cite{LeCun2015, Reichstein2019}. In hydrology and water resources, deep neural networks---including Long Short-Term Memory (LSTM) networks and Convolutional Neural Networks (CNNs)---have been extensively deployed for tasks ranging from rainfall-runoff prediction to flood inundation mapping \cite{Shen2018, Sit2020, Bentivoglio2023}. For instance, Mo et al. \cite{Mo2020} successfully applied deep learning for predicting estuarine and coastal storm surges, significantly accelerating prediction times compared to numerical simulators. 

These pure machine learning models learn complex input-output mappings directly from historical data or numerical simulation archives. Once trained, they can execute predictions in milliseconds, offering a transformative acceleration in computational speed \cite{Brunton2020, Balogun2024}. However, these models operate as ``black boxes'' that are fundamentally agnostic to the physical laws governing fluid dynamics. Consequently, they often struggle with spatial generalization, fail to ensure physical conservation (such as mass and momentum balance), and are highly susceptible to producing catastrophic, non-physical predictions when presented with out-of-distribution boundary conditions \cite{Karniadakis2021, Cuomo2022}.

\subsection{Physics-Informed Neural Networks and Gradient Pathologies}

A revolutionary resolution to the dichotomy between physical rigor and computational speed was the introduction of Physics-Informed Neural Networks (PINNs) by Raissi et al. \cite{Raissi2019}. PINNs seamlessly integrate observational data with the governing partial differential equations (PDEs) of a system. By leveraging automatic differentiation, the residuals of the PDEs (e.g., the SWE) are evaluated during the training process and added as penalty terms to the loss function. This constrains the neural network to explore only the hypothesis space that adheres to known physical laws.

Recent literature demonstrates the immense potential of PINNs for hydrodynamic applications. Sun et al. \cite{Sun2020} demonstrated surrogate modeling for fluid flows using physics-constrained learning without requiring massive simulation datasets. More specifically, studies have begun extending the PINN framework to the Shallow Water Equations. Yin et al. \cite{Yin2024} successfully solved the 1D unsteady SWE using a PINN approach, while Tian et al. \cite{Tian2025} and Dazzi et al. \cite{Dazzi2024} expanded this to 2D SWE domains, explicitly incorporating complex topographical and bathymetric source terms. 

Despite these successes, PINNs frequently suffer from severe training instabilities when applied to complex multi-physics problems. Wang et al. \cite{Wang2021} highlighted that PINNs are highly susceptible to ``gradient pathologies,'' where the gradient of the loss function with respect to the physical residuals completely dominates the gradients of the data or boundary conditions, leading to failure in convergence. This underscores the necessity for advanced optimization strategies and careful loss weighting when developing robust surrogate models.

\subsection{Graph Neural Networks and Unstructured Meshes}

A critical architectural limitation of standard PINNs is their reliance on Multi-Layer Perceptrons (MLPs) mapping continuous spatial coordinates $(x, y)$ or Convolutional Neural Networks (CNNs) operating on regular Cartesian grids. This renders them fundamentally incompatible with the unstructured, highly heterogeneous finite-volume meshes that characterize state-of-the-art hydrodynamic models like Delft3D FM. 

To bridge this structural gap, researchers have increasingly adopted Graph Neural Networks (GNNs) \cite{Scarselli2009, Gilmer2017}. GNNs treat the computational mesh naturally as a graph, where finite-volume cells act as nodes and the shared faces between them act as edges. By employing iterative message passing, GNNs learn complex spatial representations that explicitly respect the unstructured topology of the domain. In a seminal work, Pfaff et al. \cite{Pfaff2021} demonstrated that graph networks can successfully learn to simulate highly complex, mesh-based physics across fluid and solid mechanics. Building on this, recent applications by Maier et al. \cite{Maier2023} and Lino et al. \cite{Lino2023} have confirmed that GNNs can effectively emulate numerical fluxes across irregular coastal grids. Crucially, Liu \cite{Liu2026} introduced the Data-Guided Finite Volume-Informed Neural Network (FVM-PINN), emphasizing that strong-form PINNs fail on unstructured shallow water grids unless formulated with finite-volume physics and guided by sparse observational data. The fusion of GNNs with physics-informed training---yielding the Physics-Informed Graph Neural Network (PI-GNN)---has recently been proposed as the optimal architecture for flood forecasting on complex terrains, perfectly marrying geometric flexibility with physical adherence \cite{Taghizadeh2025, Balogun2025}.

\subsection{Mitigating Spectral Bias in Tidal Dynamics}

A foundational challenge in applying deep learning to estuarine hydrodynamics is the phenomenon of ``spectral bias.'' As established by Rahaman et al. \cite{Rahaman2019}, neural networks inherently prioritize learning low-frequency, smoothly varying components of a signal, whilst struggling profoundly to resolve high-frequency, highly oscillatory dynamics. 

In macrotidal estuaries, hydrodynamics are utterly dominated by high-frequency semi-diurnal and quarter-diurnal tidal oscillations. Standard neural networks tend to produce overly smoothed, inaccurate representations of these tidal waves. To counteract this, Tancik et al. \cite{Tancik2020} demonstrated that mapping low-dimensional inputs (such as time and space) into a high-dimensional periodic space using Random Fourier Feature (RFF) embeddings empowers networks to overcome spectral bias, enabling the accurate reconstruction of high-frequency functions. Applying RFF positional encoding to hydrodynamic PINNs is an emerging and highly critical requirement for accurately modelling tidal signals.

\subsection{The Gironde Estuary as a Complex Test Bed}

The Gironde estuary in south-western France presents a highly challenging and scientifically relevant environment for testing advanced hydrodynamic models. As the largest macrotidal estuary in Western Europe, it is characterized by massive tidal ranges, severe non-linear tidal asymmetry, and highly variable riverine discharges from the Garonne and Dordogne rivers \cite{Sottolichio2000, Jalón-Rojas2018}. 

Previous studies on the Gironde have extensively documented the complex interactions between tidal forcing and riverine flow, highlighting its intricate sediment transport mechanisms and cohesive sediment suspension dynamics \cite{Toublanc2015}. Because of its extreme hydrodynamic variability, traditional numerical methods are typically required to model it accurately \cite{GirondeTidal2019}. Its geometric complexity, vast physical scale, and extreme non-linear tidal dynamics make the Gironde an ideal, rigorous proving ground for the proposed PI-GNN surrogate model.

\section{Summary and Implications}

The comprehensive review of the literature reveals a clear trajectory in hydrodynamic modelling: progressing from computationally intensive physics-based solvers, through rapid but physically inconsistent data-driven models, and now converging upon physically constrained machine learning frameworks. To synthesize this progression, Table \ref{tab:lit_summary} provides a direct comparison between the current state-of-the-art surrogate modelling approaches and the methodology proposed in this thesis.

\begin{table}[htbp]
    \centering
    \caption{Summary of state-of-the-art surrogate modelling methodologies compared to the proposed Hydro-PI-GNN framework.}
    \label{tab:lit_summary}
    \vspace{0.2cm}
    \scriptsize
    \noindent\begin{tabularx}{\textwidth}{|>{\raggedright\arraybackslash\hspace{0pt}}X|>{\raggedright\arraybackslash\hspace{0pt}}X|>{\raggedright\arraybackslash\hspace{0pt}}X|>{\raggedright\arraybackslash\hspace{0pt}}X|>{\raggedright\arraybackslash\hspace{0pt}}X|>{\raggedright\arraybackslash\hspace{0pt}}X|}
    \hline
    \textbf{Study / Reference} & \textbf{Methodology} & \textbf{Application Domain} & \textbf{Mesh / Grid} & \textbf{Physical Consistency} & \textbf{Spectral Bias Addressed?} \\ \hline
    Sit et al. (2020) \cite{Sit2020} & Deep Learning (GCN) & Streamflow / Hydrology & Graph/Network & No \newline (Data-Driven) & No \\ \hline
    Mo et al. (2020) \cite{Mo2020} & Deep Learning (CNN/LSTM) & Coastal Storm Surges & Regular Grid & No \newline (Data-Driven) & No \\ \hline
    Raissi et al. (2019) \cite{Raissi2019} & PINN (Original) & General Fluid PDEs & Continuous 1D/2D & Yes (PDEs) & No \\ \hline
    Tancik et al. (2020) \cite{Tancik2020} & Fourier Features (MLP) & 1D/2D Signals & Continuous & No \newline (Data-Driven) & Yes (RFF) \\ \hline
    Yin et al. (2024) \cite{Yin2024} & Physics-Informed Neural Network (PINN) & 1D Shallow Water & Continuous 1D & Yes (SWE) & No \\ \hline
    Dazzi et al. (2024) \cite{Dazzi2024} & Advanced PINN & 2D Shallow Water & Regular Cartesian & Yes (SWE + Topo) & No \\ \hline
    Tian et al. (2025) \cite{Tian2025} & Advanced PINN & 2D Shallow Water & Regular Cartesian & Yes (SWE + Topo) & No \\ \hline
    Pfaff et al. (2021) \cite{Pfaff2021} & Graph Neural Networks (GNN) & General Fluid Mechanics & Unstructured & No (Implicit) & No \\ \hline
    Maier et al. (2023) \cite{Maier2023} & Graph Neural Networks (GNN) & Coastal Flows & Irregular / Mesh & No (Implicit) & No \\ \hline
    Liu (2026) \cite{Liu2026} & FVM-PINN & 2D Shallow Water & Finite Volume Mesh & Yes \newline (FVM-SWE) & No \\ \hline
    Taghizadeh et al. (2025) \cite{Taghizadeh2025} & PI-GNN & Flood Forecasting & Graph-based & Yes & No \\ \hline
    \textbf{This Study \newline (Hydro-PI-GNN)} & \textbf{PI-GNN with RFF Encoding} & \textbf{Macrotidal Estuary (2D SWE)} & \textbf{Unstructured \newline (Delft3D FM)} & \textbf{Yes \newline (Autograd SWE)} & \textbf{Yes \newline (Dual-Band RFF)} \\ \hline
    \end{tabularx}
\end{table}


Several critical research gaps have been identified in the current state-of-the-art:
\begin{enumerate}
    \item Purely data-driven models (e.g., \cite{Mo2020}) cannot guarantee physical consistency, limiting their operational reliability during extreme, unseen weather events.
    \item While standard PINNs (e.g., \cite{Tian2025}) have demonstrated success in solving the Shallow Water Equations, they are largely restricted to continuous coordinate mapping or regular grids, rendering them incompatible with the unstructured finite-volume meshes utilized by modern estuarine models like Delft3D FM.
    \item Standard deep learning architectures suffer from well-documented spectral bias \cite{Rahaman2019}, routinely failing to capture the high-frequency tidal oscillations that define macro-tidal estuarine dynamics without specific feature embedding \cite{Tancik2020}.
    \item There is a distinct lack of fully parametric PI-GNN surrogates capable of modelling large-scale, real-world estuaries while being explicitly conditioned on highly variable, dynamic boundary conditions (such as varying ocean tides and upstream river discharges).
\end{enumerate}

These identified gaps directly form the implications and foundational justification for this thesis. To advance the field of real-time coastal forecasting, it is necessary to move beyond simple grid-based PINNs. The literature strongly supports the development of a hybrid approach: merging the topological awareness of Graph Neural Networks \cite{Pfaff2021} with the physical rigor of autograd-based Shallow Water Equation residuals. Furthermore, the integration of Random Fourier Feature (RFF) embeddings is implied as a critical necessity to overcome the spectral bias inherent in tidal modelling. By independently developing and calibrating a Delft3D FM numerical model of the highly dynamic Gironde estuary, and subsequently training a parametric PI-GNN architecture upon its unstructured mesh, this study directly addresses the limitations of current methodologies. This approach promises to deliver a digital twin capable of unprecedented computational speed without sacrificing the physical fidelity demanded by modern coastal engineering and climate impact assessment.
